\newcommand{\figuraElementoYoungLaplace}{%
\begin{figure}[H]
\centering
\begin{tikzpicture}[>=Latex,font=\small,line cap=round,line join=round]
  \node[font=\bfseries] at (3,5.15) {(a) Vista normal};
  \fill[blue!12] (1.2,1.3) rectangle (4.8,4.0);
  \draw[very thick] (1.2,1.3) rectangle (4.8,4.0);
  \draw[<->] (1.2,0.9)--(4.8,0.9) node[midway,below] {$d\ell_1$};
  \draw[<->] (0.8,1.3)--(0.8,4.0) node[midway,left] {$d\ell_2$};
  \draw[->,red!75!black,thick] (1.2,3.15)--(0.3,3.15) node[left] {$\alpha d\ell_2$};
  \draw[->,red!75!black,thick] (4.8,3.15)--(5.7,3.15) node[right] {$\alpha d\ell_2$};
  \draw[->,green!45!black,thick] (3.0,4.0)--(3.0,4.58)
    node[midway,right] {$\alpha d\ell_1$};
  \draw[->,green!45!black,thick] (3.8,1.3)--(3.8,0.55)
    node[below] {$\alpha d\ell_1$};

  \node[font=\bfseries] at (9,4.9) {(b) Corte principal};
  \draw[very thick,blue!70!black] (6.4,2.0)
    .. controls (7.6,4.2) and (10.4,4.2) .. (11.6,2.0);
  \coordinate (L) at (7.45,3.25);
  \coordinate (R) at (10.55,3.25);
  \draw[->,ultra thick,red!75!black] (L)--++(150:1.0);
  \draw[->,ultra thick,red!75!black] (R)--++(30:1.0);
  \draw[->,ultra thick,purple!75!black] (9,3.65)--(9,2.25)
    node[midway,right] {resultante normal};
  \coordinate (C) at (9,0.65);
  \fill (C) circle (0.05) node[below] {$C_1$};
  \draw[densely dashed] (C)--(L) (C)--(R);
  \node at (9,1.25) {$R_1$};
\end{tikzpicture}
\caption{Elemento superficial alineado con las direcciones principales. Las
fuerzas tangenciales en lados opuestos producen una resultante normal asociada
con cada radio de curvatura.}
\label{fig:elemento-young-laplace}
\end{figure}%
}